\chapter{A-TP via the Exact Hadamard Factorisation: Reflected Resolvents and Tautology Comparison}
\label{v4-arith:w7-g:ATP-Hadamard}

\section{Hadamard Factorisation}
\label{v4-arith:w7-g:context}

\subsection{Setup}
\label{v4-arith:w7-g:context:obstruction}

The archimedean Tate positivity conjecture A-TP
(\Cref{v4-arith:w5-a:conj:archimedean-tate-positivity}) is the
statement that the Weil distribution \(W(f*\widetilde{f})\geq 0\) for
all \(f\) in the Paley--Wiener space \(\mathrm{PW}(\mathbb{R})\) (entire
functions of exponential type whose restriction to \(\mathbb{R}\) is
square-integrable; \(\widetilde{f}(x):=\overline{f(-x)}\)). Chapters~\ref{v4-arith:w5-j:VBstar-spectral-positivity}
and~\ref{v4-arith:w6-a:ATP-direct-obstruction} tested A-TP via de Branges
/ Hermite--Biehler spectral theory and via the digamma-kernel elementary
decomposition, respectively; each construction bottomed out at a hypothesis
equivalent to RH.

Hadamard factorisation gives
\(\xi_{\mathbb{R}}(s)=e^{a+bs}\prod_{\rho}(1-s/\rho)e^{s/\rho}\)
of \(\xi_{\mathbb{R}}\) as an entire function of order \(1\) and genus \(1\)
(\Cref{v4-arith:w5-j:prop:xi-order}). Logarithmic differentiation
yields \(\xi_{\mathbb{R}}'/\xi_{\mathbb{R}}\) as a sum over zeros plus a
linear correction. The reflected-resolvent kernel
\(\mathcal{K}_{\rho}(t)\) obtained by multiplying opposite critical-line
resolvents is not the Guinand--Weil zero summand; it is nevertheless
sign-indefinite even when \(\rho\) lies on the critical line. Thus the
Hadamard construction gives no termwise positivity theorem. The
positivity statement is the full Weil distribution, hence RH-level.

\subsection{Classical Sources}
\label{v4-arith:w7-g:sources}

The Hadamard-factorisation obstruction uses four classical sources
disjoint from the restricted-class comparison and the bounded-interval reduction.

\begin{description}
\item[Hadamard 1893.] Hadamard, J. \emph{\'Etude sur les propri\'et\'es
des fonctions enti\`eres et en particulier d'une fonction consid\'er\'ee
par Riemann.} \emph{J.~Math.~Pures Appl.}~(4) 9 (1893), 171--215.
Original Hadamard factorisation theorem for entire functions of
finite order.
\item[Guinand--Weil 1948/1952.] Guinand, A.P. \emph{A summation formula
in the theory of prime numbers.} \emph{Proc.~London Math.~Soc.}~(2) 50
(1948), 107--119. Weil, A. \emph{Sur les formules explicites de la
th\'eorie des nombres premiers} (cited above). Pair-pole
representation via Hadamard derivative; Parseval between critical-line
Fourier and zero-sum.
\item[de Bruijn 1950.] de Bruijn, N.G. \emph{The roots of trigonometric
integrals and a theorem of Hurwitz.} \emph{Duke Math.~J.}~17 (1950),
197--226. Heat-flow family \(\Xi_{\lambda}\); monotone shift of roots
towards the real axis.
\item[Newman 1976, Rodgers--Tao 2020.] Newman, C.M. \emph{Fourier
transforms with only real zeros.} \emph{Proc.~AMS} 61 (1976), 245--251.
Rodgers, B.; Tao, T. \emph{The de Bruijn--Newman constant is
non-negative.} \emph{Forum Math.~Pi} 8 (2020), e6. The de Bruijn--Newman
constant \(\Lambda\); RH iff \(\Lambda\leq 0\); unconditional lower bound
\(\Lambda\geq 0\).
\end{description}

Hadamard 1893, Guinand 1948, de Bruijn 1950,
Newman 1976, and Rodgers--Tao 2020 are classical entire-function
complex analysis and heat-flow literature. None enter the de Branges spectral comparison
de Branges / Hermite--Biehler derivation
(Chapter~\ref{v4-arith:w5-j:VBstar-spectral-positivity}); none enter
the bounded-interval reduction elementary digamma kernel analysis
(Chapter~\ref{v4-arith:w6-a:ATP-direct-obstruction}); none enter the
Mellin--Rankin--Selberg isometry of
Chapter~\ref{v4-arith:w5-a:chapter}. These sources are
disjoint from Chapters~\ref{v4-arith:closure},
\ref{v4-arith:kp-compat}, \ref{v4-arith:kms-gns-full}.

\section{The Exact Hadamard Formula for \(\xi_{\mathbb{R}}'/\xi_{\mathbb{R}}\)}
\label{v4-arith:w7-g:hadamard-formula}

\subsection{Hadamard Product Logarithmic Derivative}
\label{v4-arith:w7-g:hadamard-formula:log-derivative}

\begin{lemma}[Hadamard formula for \(\xi_{\mathbb{R}}'/\xi_{\mathbb{R}}\)]
\label{v4-arith:w7-g:lem:hadamard-log-deriv}
\ClaimStatusProvedHere. Fix the Hadamard factorisation
\eqref{v4-arith:w5-j:eq:hadamard}
\(\xi_{\mathbb{R}}(s)=e^{a+bs}\prod_{\rho}(1-s/\rho)e^{s/\rho}\) with
\(a,b\) real constants of \Cref{v4-arith:w5-j:prop:xi-order} and product
over non-trivial zeros. The logarithmic derivative satisfies
\begin{equation}
\label{v4-arith:w7-g:eq:log-deriv-Hadamard}
\frac{\xi_{\mathbb{R}}'(s)}{\xi_{\mathbb{R}}(s)}
\;=\;b+\sum_{\rho}\!\Bigl(\frac{1}{s-\rho}+\frac{1}{\rho}\Bigr),
\end{equation}
where the series converges absolutely and uniformly on compact subsets
avoiding the zero locus, after pairing \(\rho\leftrightarrow 1-\rho\).
\end{lemma}

\begin{proof}
Logarithmic differentiation of \eqref{v4-arith:w5-j:eq:hadamard}:
\(\log\xi_{\mathbb{R}}(s)=a+bs+\sum_{\rho}[\log(1-s/\rho)+s/\rho]\).
Differentiating:
\(\xi_{\mathbb{R}}'(s)/\xi_{\mathbb{R}}(s)=b+\sum_{\rho}[(-1/\rho)/(1-s/\rho)+1/\rho]
=b+\sum_{\rho}[1/(s-\rho)+1/\rho]\) after the simplification
\((-1/\rho)/(1-s/\rho)=1/(s-\rho)\). Absolute convergence: under the
pairing \(\rho\mapsto 1-\rho\), the paired term
\([1/(s-\rho)+1/(s-(1-\rho))]+[1/\rho+1/(1-\rho)]\) has both halves
absolutely summable since
\(\sum_{\rho}(|\rho(1-\rho)|)^{-1}<\infty\) by the Riemann--von Mangoldt
bound \(N(T)\asymp T\log T\).
\end{proof}

\subsection{Comparison with the Completed Zeta Formula}
\label{v4-arith:w7-g:hadamard-formula:comparison}

\begin{proposition}[partial-fraction and gamma-factor identity]
\label{v4-arith:w7-g:prop:xi-over-xi-explicit}
\ClaimStatusProvedElsewhere (Titchmarsh \emph{The Theory of the
Riemann Zeta Function} (Oxford 1986), \S3.3; Iwaniec--Kowalski \S5.9).
For \(s\notin\{0,1\}\cup\{\rho\}\),
\begin{equation}
\label{v4-arith:w7-g:eq:xi-over-xi-split}
\frac{\xi_{\mathbb{R}}'(s)}{\xi_{\mathbb{R}}(s)}
\;=\;\frac{1}{s}+\frac{1}{s-1}+\frac{1}{2}\,\psi\!\left(\frac{s}{2}\right)-\frac{\log\pi}{2}+\frac{\zeta'(s)}{\zeta(s)}.
\end{equation}
\end{proposition}

\begin{proof}
\(\xi_{\mathbb{R}}(s)=\tfrac{1}{2}s(s-1)\pi^{-s/2}\Gamma(s/2)\zeta(s)\);
logarithmic derivative of each factor:
\((s(s-1))'/(s(s-1))=1/s+1/(s-1)\),
\((\pi^{-s/2})'/\pi^{-s/2}=-\tfrac{1}{2}\log\pi\),
\((\Gamma(s/2))'/\Gamma(s/2)=\tfrac{1}{2}\psi(s/2)\), and
\(\zeta'/\zeta\).
\end{proof}

\begin{corollary}[Hadamard and explicit formula match]
\label{v4-arith:w7-g:cor:hadamard-equals-explicit}
\ClaimStatusProvedHere. Equating \eqref{v4-arith:w7-g:eq:log-deriv-Hadamard}
and \eqref{v4-arith:w7-g:eq:xi-over-xi-split}:
\begin{equation}
\label{v4-arith:w7-g:eq:hadamard-explicit-equality}
b+\sum_{\rho}\!\Bigl(\frac{1}{s-\rho}+\frac{1}{\rho}\Bigr)
\;=\;\frac{1}{s}+\frac{1}{s-1}+\frac{1}{2}\,\psi\!\left(\frac{s}{2}\right)-\frac{\log\pi}{2}+\frac{\zeta'(s)}{\zeta(s)}.
\end{equation}
This is the standard Stark identity; setting \(s=1/2+it\) and taking
the real part gives a pointwise formula for the Weil kernel.
\end{corollary}

\begin{proof}
Both sides express the same meromorphic function; comparison is
algebraic, valid away from the zero/pole locus.
\end{proof}

\begin{remark}[the two decompositions of \(\xi_{\mathbb{R}}'/\xi_{\mathbb{R}}\)]
\label{v4-arith:w7-g:rem:chapter-split}
The right-hand side of
\eqref{v4-arith:w7-g:eq:hadamard-explicit-equality} is the Weil-kernel
split of Chapter~\ref{v4-arith:w5-a:chapter}: polar terms
\(1/s+1/(s-1)\), archimedean digamma \(\tfrac{1}{2}\psi(s/2)-\tfrac{1}{2}\log\pi\)
(the kernel \(g(t)\) of Chapter~\ref{v4-arith:w6-a:ATP-direct-obstruction} after shifting argument), finite-prime
\(\zeta'/\zeta\). The left-hand side is the sum over non-trivial zeros.
The following section separates the logarithmic-derivative zero sum
from a reflected-resolvent quadratic form; the latter is not the
Guinand--Weil zero side.
\end{remark}

\section{The Reflected Pair-Pole Diagnostic}
\label{v4-arith:w7-g:pair-pole}

\subsection{Pairing Conjugate Zeros}
\label{v4-arith:w7-g:pair-pole:pairing}

\begin{definition}[pair-pole kernel]
\label{v4-arith:w7-g:def:pair-pole-kernel}
\ClaimStatusDefinitional. Write each non-trivial zero \(\rho=\beta+i\gamma\)
with \(0\leq\beta\leq 1\) (non-trivial strip), and pair with its
complex conjugate \(\bar{\rho}=\beta-i\gamma\) and functional-equation
partner \(1-\rho=1-\beta-i\gamma\) (and \(1-\bar{\rho}=1-\beta+i\gamma\)).
For \(s=1/2+it\) on the critical line, define the \emph{pair-pole
kernel}
\begin{equation}
\label{v4-arith:w7-g:eq:pair-pole-kernel}
\mathcal{K}_{\rho}(t)\;:=\;\mathrm{Re}\!\left[\frac{1}{(1/2+it-\rho)(1/2-it-\rho)}\right].
\end{equation}
When \(\rho\) is paired with the reflected pole in the \(t\)-variable,
the denominator becomes
\(((1/2-\beta)+i(t-\gamma))((1/2-\beta)-i(t+\gamma))\), explicit
sign structure computed below.
\end{definition}

\begin{lemma}[explicit form of the pair-pole kernel]
\label{v4-arith:w7-g:lem:pair-pole-explicit}
\ClaimStatusProvedHere. For \(\rho=\beta+i\gamma\) with
\(\beta\in[0,1]\), \(\gamma\in\mathbb{R}\), and \(\alpha:=1/2-\beta\),
\begin{equation}
\label{v4-arith:w7-g:eq:K-explicit}
\mathcal{K}_{\rho}(t)
\;=\;\frac{\alpha^{2}+t^{2}-\gamma^{2}}{[\alpha^{2}+(t-\gamma)^{2}][\alpha^{2}+(t+\gamma)^{2}]}.
\end{equation}
When \(\beta=1/2\) (RH), this simplifies to
\(\mathcal{K}_{\rho}(t)\big|_{\beta=1/2}=1/(t^{2}-\gamma^{2})\),
which is a sign-indefinite distribution: positive for \(|t|>|\gamma|\),
negative for \(|t|<|\gamma|\), with distributional poles at \(t=\pm\gamma\).
The pair-pole sum \(\int|\widehat{f}|^{2}\mathcal{K}_{\rho}\,dt\) is
a principal-value term; it is not sign-definite term by term
(Lemma~\ref{v4-arith:w7-g:lem:residue-extract}).
\end{lemma}

\begin{proof}
Write \(\alpha:=1/2-\beta\) and compute:
\((1/2+it-\rho)(1/2-it-\rho)=(\alpha+i(t-\gamma))(\alpha-i(t+\gamma))\).
Expanding: \(\alpha^{2}-2i\alpha\gamma+(t^{2}-\gamma^{2})\). The real part of
the reciprocal is
\([\alpha^{2}+(t^{2}-\gamma^{2})]/[(\alpha^{2}+(t^{2}-\gamma^{2}))^{2}+4\alpha^{2}\gamma^{2}]\),
and the denominator equals \([\alpha^{2}+(t-\gamma)^{2}][\alpha^{2}+(t+\gamma)^{2}]\)
by expanding both sides. The specialisation \(\alpha=0\)
yields \(\mathcal{K}_{\rho}(t)\big|_{\alpha=0}=(t^{2}-\gamma^{2})/(t^{2}-\gamma^{2})^{2}=1/(t^{2}-\gamma^{2})\);
the symmetric contour gives the principal-value distribution of
\Cref{v4-arith:w7-g:lem:residue-extract}.
\end{proof}

\begin{remark}[the sign of \(\mathcal{K}_{\rho}(t)\)]
\label{v4-arith:w7-g:rem:off-critical-sign}
For any \(\beta\in[0,1]\), the numerator \(\alpha^{2}+t^{2}-\gamma^{2}\)
of \eqref{v4-arith:w7-g:eq:K-explicit} is negative for
\(|t|<\sqrt{\gamma^{2}-\alpha^{2}}\) (when \(|\gamma|>|\alpha|\)), while
the denominator is always positive. Hence \(\mathcal{K}_{\rho}(t)<0\)
on this interval, including in the on-critical case \(\alpha=0\) where
\(\mathcal{K}_{\rho}(t)<0\) for \(|t|<|\gamma|\). Thus no
term-by-term positivity follows from the pair-pole kernel.
\end{remark}

\subsection{Residue Extraction and Dirac Limits}
\label{v4-arith:w7-g:pair-pole:residue}

\begin{lemma}[principal value at an on-critical pole]
\label{v4-arith:w7-g:lem:residue-extract}
\ClaimStatusProvedHere. Let \(\rho=1/2+i\gamma\) be on the critical
line (\(\beta=1/2\), \(\gamma\neq 0\), \(\gamma\in\mathbb{R}\)). For \(f\in\mathrm{PW}(\mathbb{R})\)
real and even with \(\widehat{f}\) smooth at \(t=\pm\gamma\),
\begin{equation}
\label{v4-arith:w7-g:eq:residue-on-critical}
\lim_{\varepsilon\downarrow 0}\,\mathrm{Re}\!\int_{\mathbb{R}}|\widehat{f}(t)|^{2}\cdot\frac{dt}{(1/2+\varepsilon+it-\rho)(1/2+\varepsilon-it-\rho)}
\;=\;\mathrm{PV}\!\int_{\mathbb R}\frac{|\widehat f(t)|^{2}}{t^{2}-\gamma^{2}}\,dt,
\end{equation}
with the reflected pole convention of
\eqref{v4-arith:w7-g:eq:pair-pole-kernel}. This principal value is not
sign-definite.
\end{lemma}

\begin{proof}
With the reflected pole convention the regularised real kernel is
\[
\mathrm{Re}\frac{1}{(\varepsilon+i(t-\gamma))(\varepsilon-i(t+\gamma))}
=\frac{\varepsilon^{2}+t^{2}-\gamma^{2}}
{(\varepsilon^{2}+(t-\gamma)^{2})(\varepsilon^{2}+(t+\gamma)^{2})}.
\]
As \(\varepsilon\downarrow0\) this converges in distributions to
\(\mathrm{PV}(t^{2}-\gamma^{2})^{-1}\). Pairing with the smooth
function \(|\widehat f|^{2}\) gives the displayed identity.
\end{proof}

\begin{theorem}[the reflected pair-pole kernel is not the Guinand--Weil summand]
\label{v4-arith:w7-g:thm:guinand-weil-pair-pole}
\ClaimStatusProvedHere. For \(f\in\mathrm{PW}(\mathbb{R})\) real and
even, the reflected-resolvent expression
\begin{equation}
\label{v4-arith:w7-g:eq:guinand-weil}
\mathcal R_{\mathrm{ref}}(f)\;:=\;\lim_{T\to\infty}\sum_{|\mathrm{Im}\rho|\leq T}R_{\rho}[f],
\qquad
R_{\rho}[f]\;:=\;\mathrm{PV}\!\int_{\mathbb{R}}|\widehat{f}(t)|^{2}\,\mathcal{K}_{\rho}(t)\,dt,
\end{equation}
where \(\mathcal{K}_{\rho}\) is the reflected pair-pole kernel
\eqref{v4-arith:w7-g:eq:K-explicit}, the sum is over non-trivial
zeros of \(\xi_{\mathbb{R}}\), each paired once with its conjugate. When
the symmetric limit exists, it is not the zero side of the
Guinand--Weil formula. It is a diagnostic quadratic form obtained by
multiplying reflected resolvents.
\end{theorem}

\begin{proof}
It is enough to compare a single on-critical conjugate pair.  For
\(\rho=1/2+i\gamma\),
\[
\frac{1}{1/2+it-\rho}+\frac{1}{1/2+it-\bar\rho}
=-i\left(\frac{1}{t-\gamma}+\frac{1}{t+\gamma}\right)
=-\frac{2it}{t^{2}-\gamma^{2}} .
\]
Thus the zero contribution to the imaginary part of the logarithmic
derivative is the odd principal-value distribution
\(-2t\,\mathrm{PV}(t^{2}-\gamma^{2})^{-1}\).  The reflected kernel
\(\mathcal K_{\rho}(t)=\mathrm{PV}(t^{2}-\gamma^{2})^{-1}\) is even.
It therefore cannot be the residue summand obtained from
\(\xi_{\mathbb R}'/\xi_{\mathbb R}\). The displayed
\(\mathcal R_{\mathrm{ref}}\) is consequently not a Guinand--Weil
decomposition.
\end{proof}

\subsection{Under RH: No Termwise Positivity}
\label{v4-arith:w7-g:pair-pole:RH-positivity}

\begin{proposition}[the on-critical pair-pole term is sign-indefinite]
\label{v4-arith:w7-g:prop:pointwise-RH}
\ClaimStatusProvedHere. Let \(\rho=1/2+i\gamma\) be on the critical
line with \(\gamma\neq0\). Then
\begin{equation}
\label{v4-arith:w7-g:eq:pointwise-RH}
\mathcal K_{\rho}(t)=\mathrm{PV}\frac{1}{t^{2}-\gamma^{2}}.
\end{equation}
It is sign-indefinite. In particular there are real even
\(f\in\mathrm{PW}(\mathbb R)\) for which the single principal-value
term \(R_{\rho}[f]\) is negative. RH gives positivity of the full Weil
distribution, not positivity of each reflected pair-pole term.
\end{proposition}

\begin{proof}
This is \Cref{v4-arith:w7-g:lem:residue-extract}. On
\(|t|<|\gamma|\) the kernel is negative, away from the poles. Choose
\(\widehat f\) real, even, smooth, and supported in a sufficiently
small neighbourhood of \(0\). Then the principal-value integral is an
ordinary integral over a region where the kernel is negative.
\end{proof}

\section{Tautology Comparison: Where Positivity Becomes Circular}
\label{v4-arith:w7-g:tautology-comparison}

Even under RH, a single pair-pole kernel is sign-indefinite. The
following proposition records the failure of the pointwise construction.

\begin{proposition}[pointwise pair-pole positivity fails]
\label{v4-arith:w7-g:prop:circular-reduction}
\ClaimStatusProvedHere. The statement
``\(\mathcal{K}_{\rho}(t)\geq 0\) pointwise for every non-trivial zero
\(\rho\)'' is false for the zeta zero set. Hence any attempt to prove
A-TP by \emph{term-by-term} pointwise non-negativity of
\(\mathcal{K}_{\rho}(t)\) is not merely circular; it has no valid
on-critical case.
\end{proposition}

\begin{proof}
For an on-critical zero, \(\alpha=0\), and
\(\mathcal K_{\rho}(0)=-\gamma^{-2}<0\). Thus RH itself does not make
the pair-pole kernel pointwise non-negative.
\end{proof}

\begin{remark}[pointwise positivity versus integrated positivity]
\label{v4-arith:w7-g:rem:ruled-out-pointwise}
Proposition~\ref{v4-arith:w7-g:prop:circular-reduction} establishes
that pointwise \(\mathcal{K}_{\rho}(t)\geq 0\) cannot be the mechanism
for A-TP. Sum-level non-negativity belongs to the actual Weil
distribution; the reflected-resolvent form
\(\mathcal R_{\mathrm{ref}}\) is not that distribution. The
Hadamard logarithmic derivative has the same RH ceiling as the de Branges construction
(Chapter~\ref{v4-arith:w5-j:VBstar-spectral-positivity}) and the
digamma-kernel construction (Chapter~\ref{v4-arith:w6-a:ATP-direct-obstruction}).
\end{remark}

\begin{corollary}[the valid reformulations coincide]
\label{v4-arith:w7-g:cor:three-constructions-equivalent}
\ClaimStatusProvedHere. The two valid positivity statements
\begin{enumerate}
\item \emph{(Weil)} \(W(f*\widetilde{f})\geq 0\) for all
\(f\in\mathrm{PW}(\mathbb{R})\);
\item \emph{(Hermite--Biehler)} the de Branges datum
\(E_{\xi}(z)=A_{\xi}(z)-iB_{\xi}(z)\) (defined in
Chapter~\ref{v4-arith:w5-j:VBstar-spectral-positivity}) is
Hermite--Biehler and satisfies
\(\mathrm{div}(A_{\xi})=\mathrm{div}(\Xi)\);
\end{enumerate}
are equivalent and each equivalent to RH. The pointwise pair-pole
positivity statement is false and is not a third reformulation.
\end{corollary}

\begin{proof}
(1) \(\Leftrightarrow\) RH: \Cref{v4-arith:w5-a:cor:weil-positivity-RH}.
(2) \(\Leftrightarrow\) RH: Lagarias' Hermite--Biehler theorem and
\Cref{v4-arith:w5-j:thm:zeros-of-Xi-under-HB}.
\Cref{v4-arith:w7-g:prop:circular-reduction} excludes the pointwise
pair-pole statement.
\end{proof}

\section{de Bruijn--Newman \(\Lambda\leq 0\) is RH}
\label{v4-arith:w7-g:de-Bruijn-Newman}

The pair-pole pointwise positivity statement is false. The
de Bruijn--Newman bound \(\Lambda\leq 0\) is a valid construction to A-TP
only because it is RH itself.

\subsection{The de Bruijn--Newman Family}
\label{v4-arith:w7-g:de-Bruijn-Newman:family}

\begin{definition}[de Bruijn--Newman heat flow]
\label{v4-arith:w7-g:def:dBN-family}
\ClaimStatusDefinitional (de Bruijn 1950; Newman 1976). Let
\(\Xi(t):=\xi_{\mathbb{R}}(1/2+it)\). Its Fourier representation
\(\Xi(t)=\int_{0}^{\infty}\Phi(u)\cos(tu)\,du\) with
\begin{equation}
\label{v4-arith:w7-g:eq:Phi}
\Phi(u)\;:=\;\sum_{n=1}^{\infty}\pi n^{2}(2\pi n^{2}e^{4u}-3)e^{5u/2}\exp(-\pi n^{2}e^{4u})
\end{equation}
(Polya 1927, Titchmarsh \S10.1). Define the one-parameter family
\(\Xi_{\lambda}(t):=\int_{0}^{\infty}\Phi(u)e^{\lambda u^{2}}\cos(tu)\,du\)
for \(\lambda\in\mathbb{R}\). The \emph{de Bruijn--Newman constant}
\(\Lambda\) is the infimum of \(\lambda\) for which all zeros of
\(\Xi_{\lambda}\) are real.
\end{definition}

\begin{theorem}[de Bruijn--Newman bounds]
\label{v4-arith:w7-g:thm:dBN-bounds}
\ClaimStatusProvedElsewhere (Newman 1976; Rodgers--Tao 2020).
\begin{enumerate}
\item \(\Lambda\leq 1/2\) (de Bruijn 1950: for \(\lambda\geq 1/2\),
\(\Xi_{\lambda}\) has only real zeros; Newman 1976 introduced
\(\Lambda\) as an infimum).
\item \(\Lambda\geq 0\) (Rodgers--Tao 2020, unconditional).
\item RH \(\Leftrightarrow\) \(\Lambda\leq 0\) (Riemann 1859 for the direction
\(\Lambda\leq 0\Rightarrow\) all zeros on critical line; Newman 1976 for
the equivalence statement).
\item Newman (1976) conjectured \(\Lambda=0\); Rodgers--Tao 2020 established
\(\Lambda\geq 0\), so \(\Lambda=0\) is the combined best bound.
\end{enumerate}
\end{theorem}

\begin{proof}
Newman 1976 \emph{Proc.~AMS} 61; de Bruijn 1950 \emph{Duke Math.~J.}~17;
Rodgers--Tao 2020 \emph{Forum Math.~Pi} 8; and the classical
equivalence RH \(\Leftrightarrow\) \(\Xi\) has only real zeros
(Riemann 1859).
\end{proof}

\subsection{\(\Lambda\leq 0\) Implies A-TP via RH}
\label{v4-arith:w7-g:de-Bruijn-Newman:sufficiency}

\begin{lemma}[critical-line reformulation of \(\Lambda\leq 0\)]
\label{v4-arith:w7-g:lem:Lambda-leq-zero}
\ClaimStatusProvedHere. The statement \(\Lambda\leq 0\) is equivalent to:
every non-trivial zero of \(\xi_{\mathbb{R}}\) lies on the critical line
\(\mathrm{Re}\,s=1/2\).
\end{lemma}

\begin{proof}
\(\Lambda\leq 0\Leftrightarrow\Xi_{0}=\Xi\) has only real zeros
\(\Leftrightarrow\xi_{\mathbb{R}}(1/2+it)=0\) for \(t\) real only
\(\Leftrightarrow\) every non-trivial zero \(\rho\) of \(\xi_{\mathbb{R}}\)
satisfies \(\mathrm{Re}\,\rho=1/2\). This is RH.
\end{proof}

\begin{corollary}[\(\Lambda\leq 0\) implies A-TP]
\label{v4-arith:w7-g:cor:Lambda-implies-ATP}
\ClaimStatusProvedHereConditional. Under \(\Lambda\leq 0\) (equivalent to RH),
A-TP holds.
\end{corollary}

\begin{proof}
Lemma~\ref{v4-arith:w7-g:lem:Lambda-leq-zero} reduces \(\Lambda\leq 0\) to
RH. Then \Cref{v4-arith:w5-a:cor:weil-positivity-RH} gives A-TP.
\end{proof}

\begin{remark}[\(\Lambda\leq 0\) is equivalent to RH]
\label{v4-arith:w7-g:rem:Lambda-not-new}
\(\Lambda\leq 0\) is equivalent to RH, not weaker. The
Rodgers--Tao theorem \(\Lambda\geq 0\) places the constant at the
Newman barrier. change
on \(\Lambda\) is change towards RH, not a weaker substitute. The
Hadamard-factorisation construction is tautological: it reduces
A-TP to RH without passing through any strictly weaker intermediate.

The candidate ``\(\Lambda\leq\lambda_{0}\) for some \(\lambda_{0}>0\)
still implies A-TP'' is not supplied by the de Bruijn--Newman theory.
The heat flow parameterises the zero-reality problem; the value
\(\lambda=0\) is exactly RH.
\end{remark}

\begin{proposition}[zero-free regions do not give a pair-pole positivity criterion]
\label{v4-arith:w7-g:prop:weaker-condition-elementary}
\ClaimStatusProvedHere. A classical zero-free region, including the
Korobov--Vinogradov region, does not imply A-TP or a partial-support
A-TP statement through the pair-pole kernels. It excludes zeros near
the edge of the critical strip; it does not make the central
principal-value kernels positive.
\end{proposition}

\begin{proof}
\Cref{v4-arith:w7-g:prop:pointwise-RH} shows that even an on-critical
zero gives a sign-indefinite principal-value kernel. Moving zeros
away from the edge of the strip cannot turn these central kernels
into positive kernels. A zero-free region is a zero-location theorem;
A-TP is a positivity theorem for the full Weil sum.
\end{proof}

\begin{remark}[zero-free region boundary]
\label{v4-arith:w7-g:rem:zero-free-region-standard}
The Korobov--Vinogradov region is an edge exclusion near
\(\mathrm{Re}\,s=1\), with the reflected exclusion near
\(\mathrm{Re}\,s=0\). It is not a central critical-line positivity
statement. No partial-support A-TP closure is recorded here from that
input.
\end{remark}

\section{Zero Density and Pair-Pole Convergence}
\label{v4-arith:w7-g:zero-density}

\subsection{Convergence of the Pair-Pole Sum}
\label{v4-arith:w7-g:zero-density:convergence}

\begin{proposition}[reflected pair-pole form converges distributionally after pairing]
\label{v4-arith:w7-g:prop:pair-pole-converges}
\ClaimStatusProvedHere. For \(f\in\mathrm{PW}(\mathbb{R})\), the
reflected pair-pole form \eqref{v4-arith:w7-g:eq:guinand-weil} is a paired
principal-value distribution. It is not, in general, an absolutely
convergent sum of ordinary integrals. Absolute convergence holds only
after imposing vanishing near the finitely many zero ordinates met by
\(\mathrm{supp}\,\widehat f\), or after keeping the explicit
regularisation.
\end{proposition}

\begin{proof}
Away from \(t=\pm\gamma\), the displayed kernel has the stated
\(|\gamma|^{-2}\) decay after pairing, and the zero-counting estimate
controls the tail. At \(t=\pm\gamma\), however, the on-critical kernel
is \(\mathrm{PV}(t^{2}-\gamma^{2})^{-1}\), whose absolute value is not
locally integrable. Thus the reflected-resolvent expression is a
principal-value distribution unless the test function kills those
poles or a regularisation is retained.
\end{proof}

\begin{remark}[sum-level positivity versus pointwise positivity]
\label{v4-arith:w7-g:rem:sum-vs-pointwise}
For \(\rho=\beta+i\gamma\) with \(\alpha=1/2-\beta\), the functional-equation
partner \(1-\rho=(1-\beta)-i\gamma\) has \(\alpha'=1/2-(1-\beta)=-\alpha\).
Since \eqref{v4-arith:w7-g:eq:K-explicit} depends on \(\alpha^{2}\),
\(\mathcal{K}_{1-\rho}(t)=\mathcal{K}_{\rho}(t)\); the two kernels are equal.
The sum-level non-negativity of the Weil distribution is equivalent to
RH by \Cref{v4-arith:w5-a:cor:weil-positivity-RH}. It is not the
pointwise positivity of \(\mathcal K_\rho\), and it is not the
positivity of the reflected-resolvent form
\(\mathcal R_{\mathrm{ref}}\).
\end{remark}

\subsection{Exact Domain of the Hadamard-Factorisation Construction}
\label{v4-arith:w7-g:zero-density:form}

\begin{theorem}[Hadamard-factorisation construction: exact domain]
\label{v4-arith:w7-g:thm:honest-form}
\ClaimStatusProvedHere. The exact Hadamard-factorisation obstruction on
A-TP has the following outcomes.
\begin{enumerate}
\item \emph{Reflected pair-pole diagnostic}
(\Cref{v4-arith:w7-g:thm:guinand-weil-pair-pole}): unconditional;
explicit pointwise form \eqref{v4-arith:w7-g:eq:K-explicit}, but not
the Guinand--Weil zero side.
\item \emph{No termwise positivity under RH}
(Proposition~\ref{v4-arith:w7-g:prop:pointwise-RH}): under RH,
\(\mathcal{K}_{\rho}(t)=\mathrm{PV}(t^{2}-\gamma^{2})^{-1}\), and
single pair-pole terms can be negative.
\item \emph{Valid reformulations coincide}
(\Cref{v4-arith:w7-g:cor:three-constructions-equivalent}): Weil positivity
and Hermite--Biehler positivity are equivalent to RH. Pointwise
pair-pole positivity is false
(Proposition~\ref{v4-arith:w7-g:prop:circular-reduction}).
\item \emph{de Bruijn--Newman sufficient condition}
(\Cref{v4-arith:w7-g:cor:Lambda-implies-ATP}): \(\Lambda\leq 0\)
implies A-TP; but \(\Lambda\leq 0\Leftrightarrow\) RH, so this is not
weaker than RH.
\item \emph{Zero-free region boundary}
(\Cref{v4-arith:w7-g:prop:weaker-condition-elementary}): the
Korobov--Vinogradov-style zero-free region does not give a
pair-pole positivity criterion or a partial-support A-TP closure.
\end{enumerate}
The genuine content of the Hadamard-factorisation construction is the
logarithmic-derivative identity and the reflected-resolvent diagnostic.
It does not give a termwise Weil decomposition and does not close A-TP
below the RH threshold.
\end{theorem}

\begin{proof}
Combine the cited results.
\end{proof}

\begin{remark}[the tautology map]
\label{v4-arith:w7-g:rem:tautology-map}
The three constructions to A-TP considered across the restricted-class comparison, the bounded-interval reduction, and
the analytic refinement are:
\begin{description}
\item[the de Branges spectral comparison (de Branges).] HB\(^{+}\) + BD\(^{+}\) + Det + SH;
tautology level: each hypothesis is RH-equivalent or strictly finer.
\item[the bounded-interval reduction (digamma).] Bounded-interval positivity on
\([-t_{\ast},t_{\ast}]\) after elementary reductions, with the
finite-prime and polar functional \(R[f]\) retained; tautology level:
the full inequality is RH-equivalent.
\item[the Hadamard reduction (Hadamard).] Logarithmic derivative and reflected pair-pole diagnostic;
failure mode: pointwise positivity is false even on the critical line
(\Cref{v4-arith:w7-g:prop:circular-reduction}).
\end{description}
All three constructions bottom out at RH. The content each provides is a
\emph{different local rewriting} of the same positivity question. The
Hadamard construction offers the pointwise reflected-resolvent formula
\eqref{v4-arith:w7-g:eq:K-explicit}, which clarifies the sign
obstruction but does not give the Weil zero side.
\end{remark}

\section{Conclusion}
\label{v4-arith:w7-g:conclusion}

The Hadamard factorisation
\(\xi_{\mathbb{R}}(s)=e^{a+bs}\prod_{\rho}(1-s/\rho)e^{s/\rho}\)
gives the logarithmic derivative
\[
\frac{\xi_{\mathbb R}'(s)}{\xi_{\mathbb R}(s)}
=b+\sum_{\rho}\left(\frac{1}{s-\rho}+\frac{1}{\rho}\right).
\]
The reflected kernel \eqref{v4-arith:w7-g:eq:K-explicit} is not the
Guinand--Weil zero summand. Pointwise non-negativity of
\(\mathcal{K}_{\rho}\) fails even under RH
(Proposition~\ref{v4-arith:w7-g:prop:circular-reduction}), so the
Hadamard construction cannot prove A-TP by termwise positivity. The
de Bruijn--Newman constant \(\Lambda\leq 0\) is likewise equivalent to
RH (Lemma~\ref{v4-arith:w7-g:lem:Lambda-leq-zero}). The
Korobov--Vinogradov zero-free region gives no pair-pole positivity
criterion. The valid reformulations of A-TP remain Weil positivity
and the Hermite--Biehler condition
(Corollary~\ref{v4-arith:w7-g:cor:three-constructions-equivalent}).

\begin{remark}[the non-tautological content of the Hadamard construction]
\label{v4-arith:w7-g:rem:consequence}
The Hadamard construction yields the exact logarithmic derivative of
\(\xi_{\mathbb R}\) and an explicit reflected-resolvent kernel
\eqref{v4-arith:w7-g:eq:K-explicit}. Pointwise pair-pole positivity is false
(Proposition~\ref{v4-arith:w7-g:prop:circular-reduction}), so the
construction does not close A-TP via this mechanism. Classical zero-free
regions remain zero-location input, not positivity input.
\end{remark}
